\documentclass[12pt,a4paper]{article}
\usepackage[utf8]{inputenc}
\usepackage[T1]{fontenc}
\usepackage{amsmath,amssymb}
\usepackage{enumitem}
\usepackage{geometry}
\usepackage{array}
\usepackage{multicol}
\geometry{margin=2.2cm}
\setlist{itemsep=0.25em, topsep=0.3em}
\title{Aula 13 -- Técnicas de Controle de Concorrência}
\author{Banco de Dados -- Guia de Revisão}
\date{}
\begin{document}
\maketitle
\tableofcontents
\newpage

\section{Objetivo da aula}
Esta aula mostra como o SGBD controla transações concorrentes usando bloqueios. O objetivo é preservar isolamento e consistência sem obrigar todas as transações a executarem serialmente.

\section{Por que controlar concorrência?}
Transações simultâneas podem acessar os mesmos dados. Sem controle, podem ocorrer leitura suja, atualização perdida, inconsistência e planos não serializáveis.

\section{Conflitos}
Conflitos importantes:
\begin{itemize}
    \item \textbf{read-write}: uma transação lê e outra escreve o mesmo item.
    \item \textbf{write-read}: uma transação escreve e outra lê o mesmo item.
    \item \textbf{write-write}: duas transações escrevem o mesmo item.
\end{itemize}
Leitura com leitura não conflita.

\section{Bloqueio binário}
No bloqueio binário, um item está bloqueado ou desbloqueado. Problema: não diferencia leitura de escrita, então reduz concorrência.

\section{Bloqueio compartilhado e exclusivo}
\begin{itemize}
    \item \textbf{Compartilhado} (S): usado para leitura.
    \item \textbf{Exclusivo} (X): usado para escrita.
\end{itemize}

\section{Matriz de compatibilidade}
\begin{center}
\begin{tabular}{|c|c|c|}
\hline
 & S & X \\
\hline
S & compatível & incompatível \\
X & incompatível & incompatível \\
\hline
\end{tabular}
\end{center}

Consequência:
\begin{itemize}
    \item Muitos leitores podem ler juntos.
    \item Escrita exige exclusividade.
\end{itemize}

\section{Upgrade e downgrade}
\begin{itemize}
    \item \textbf{Upgrade}: S para X. Ex.: a transação lia e agora precisa escrever.
    \item \textbf{Downgrade}: X para S. Ex.: não precisa mais escrever, só manter leitura.
\end{itemize}

\section{2PL -- bloqueio em duas fases}
O protocolo 2PL exige que toda transação tenha duas fases:
\begin{enumerate}
    \item \textbf{Fase de expansão}: pode adquirir bloqueios, mas não liberar.
    \item \textbf{Fase de encolhimento}: pode liberar bloqueios, mas não adquirir novos.
\end{enumerate}

Regra essencial: depois do primeiro \texttt{unlock}, não pode haver novo \texttt{lock}.

\subsection{Exemplo que viola 2PL}
\[
lock(A),\ read(A),\ unlock(A),\ lock(B),\ read(B),\ unlock(B)
\]
Viola porque adquire \texttt{lock(B)} depois de liberar \texttt{A}.

\subsection{Exemplo que respeita 2PL}
\[
lock(A),\ lock(B),\ read(A),\ write(B),\ unlock(A),\ unlock(B)
\]
Todos os bloqueios foram adquiridos antes do primeiro unlock.

\section{O que 2PL garante?}
O 2PL garante serializabilidade por conflito. Mas o 2PL básico pode gerar deadlock.

\section{2PL conservador}
Também chamado de 2PL estático. A transação obtém todos os bloqueios antes de iniciar.

\textbf{Vantagem}: previne deadlock.\newline
\textbf{Desvantagem}: reduz concorrência, pois prende recursos antes de usar.

\section{2PL estrito}
Mantém bloqueios exclusivos até o fim da transação. Ajuda a evitar abortos em cascata e facilita recuperação.

\section{2PL rigoroso}
Mantém bloqueios compartilhados e exclusivos até o fim. É mais restritivo que o 2PL estrito.

\section{Deadlock}
Deadlock ocorre quando transações esperam circularmente.

Exemplo:
\begin{itemize}
    \item $T_1$ bloqueia $A$ e espera $B$.
    \item $T_2$ bloqueia $B$ e espera $A$.
\end{itemize}
Nenhuma progride sem intervenção.

Soluções comuns: detectar e abortar uma transação; prevenir por protocolo; usar timeout.

\section{Starvation}
Starvation ocorre quando uma transação espera indefinidamente porque sempre é preterida. Política justa de fila pode reduzir o problema.

\section{Pegadinhas comuns}
\begin{itemize}
    \item 2PL garante serializabilidade, mas não garante ausência de deadlock.
    \item 2PL conservador evita deadlock, mas reduz concorrência.
    \item 2PL estrito mantém bloqueios exclusivos até o fim.
    \item 2PL rigoroso mantém todos os bloqueios até o fim.
\end{itemize}

\section{Desafios}
\subsection*{Desafio 1}
$T_1$ tem bloqueio compartilhado em $X$. $T_2$ quer escrever em $X$. O que acontece?

\subsection*{Desafio 2}
A sequência abaixo respeita 2PL?
\[
lock(A),\ read(A),\ lock(B),\ write(B),\ unlock(B),\ unlock(A)
\]

\section{Resolução dos desafios}
\textbf{Desafio 1}: $T_2$ precisa de bloqueio exclusivo, incompatível com S. Deve esperar ou ser abortada pelo protocolo.

\textbf{Desafio 2}: sim. Todos os locks ocorrem antes do primeiro unlock.

\section{Checklist final}
\begin{itemize}
    \item Sei a matriz S/X?
    \item Sei identificar upgrade e downgrade?
    \item Sei dizer se uma sequência viola 2PL?
    \item Sei diferenciar 2PL básico, conservador, estrito e rigoroso?
    \item Sei reconhecer deadlock e starvation?
\end{itemize}

\end{document}
